\color{uniblue}\section[Выбор уставок устройства резервирования при отказе выключателя]{\sloppy Выбор уставок устройства резервирования при отказе \mbox{выключателя}}
\color{black}

\begin{enumerate}[label=\arabic{section}.\arabic*, labelsep=4pt, leftmargin=0pt, itemindent=57pt]

\item
Ток срабатывания максимального измерительного органа тока УРОВ выбирается по условию обеспечения чувствительности:

\begin{equation}
\label{eq:urov1}
%I_{\textsl{ср.}} \leq \frac{I_{\textsl{мин}}}{k_{\textsl{ч}}} ,
I_{\textsl{ср.}} = \frac{I_{\textsl{мин}}}{k_{\textsl{ч}}} ,
\end{equation}
\begin{eqexpl}[15mm]
\item{$I_{\textsl{мин}}$} минимальный ток защищаемого объекта в аварийном режиме, А;
\item{$k_{\textsl{ч}}$} требуемый коэффициент чувствительности.
\end{eqexpl}
\item \label{urov:punkt}
Ток срабатывания рекомендуется выбирать минимальным из диапазона (5-10)\% от $I_{\textsl{ном}}$, где $I_{\textsl{ном}}$ - номинальный ток защищаемого объекта.

\item
Время действия УРОВ должно быть больше времени действия защиты на отключение:
\begin{equation}
\label{eq:urov2}
t_{\textsl{УРОВ}} = t_{\textsl{откл.выкл.}}+t_{\textsl{воз.РЗ}}+t_{\textsl{ош.РВ}}+{\Delta t},
\end{equation}
\begin{eqexpl}[15mm]
\item{$t_{\textsl{откл.выкл.}}$} время отключение выключателя (для современных выключателей составляет 0,05 – 0,1 с), с учетом времени срабатывания выходных реле устройства (0,01 с);
\item{$t_{\textsl{воз.РЗ}}$} время возврата реле защиты (для микропроцессорных терминалов составляет 0,025 с);
\item{$t_{\textsl{ош.РВ}}$} время допустимой погрешности реле времени УРОВ (для микропроцессорных терминалов составляет 0,025 с);
\item{${\Delta t}$}  запас по времени (как правило 0,1 с).
\end{eqexpl}

\end{enumerate}